\SourceLeaf{287}%
Wert $360^0$ umso näher kommt, je kleiner $\omega, \omega'$ sind. D.\,h. die Tangente an die Kurve hat ihre Neigung vom Anfang bis zum Ende um $360^0$ geändert.

Schneidet sich die Kurve einmal selbst, so ändern sich diese Schlüsse. Sind nämlich die Endpunkte von $s = \sigma$, $s' = \sigma'$ identisch, so ändert sich $(s', s)$ bei diesen Werten sprunghaft. Hält man nämlich $s = \sigma$ konstant, so ändert sich $(s', s)$, wenn $s'$ den Wert $\sigma'$ erreicht, auf einmal um $\pm 180^0$. Um zu entscheiden, welcher Wert hier gilt, muss man die Werte für $s = \sigma \mp \omega$ betrachten, wobei $\omega$ unendlich klein ist. Man sieht, dass, wenn $s' - \sigma'$ von einem [negativen] Wert [an] immer kleiner werdende negative Werte annimmt und durch $0$ wieder wächst, $(s', \sigma - \omega)$ um eine Größe zu- oder abnimmt, die sich nur unendlich wenig von $180^0$ unterscheidet, je nachdem, ob man bei $s'$ von der inneren auf die äußere Seite der Linie $s'$ übergeht oder umgekehrt. Für $(s', \sigma + \omega)$ ist es dagegen umgekehrt. Man sieht also, dass für ein konstantes $s'$, das größer als $\sigma'$ ist, $(s', s)$ beim Durchgang durch den Wert $s = \sigma$ im Allgemeinen auf einmal um $360^0$ ab(zu)nimmt und dass die Änderung von $(s', \sigma)$, wenn $s'$ den Wert $\sigma'$ durchläuft, um $\pm 180^0$ eigentlich mit beiden Vorzeichen zugleich gilt, je nachdem, ob man $s = \sigma$ als den letzten der \textls{bis} $\sigma$ oder als den ersten der \textls{von} $\sigma$ an weiter zu betrachtenden Werte auffasst. Der Körper hat hier die Gestalt einer Wendeltreppe. Die Änderung der Richtung der Tangente beträgt nun entweder $720^0$ oder $0^0$. Diese Betrachtung erübrigt sich aber ganz, wenn man die vorliegende Methode auf den Fall beschränkt, in dem sich die Kurve \textls{nicht} selbst schneidet.

\begin{center}
[8.]
\end{center}

Überhaupt lässt sich in dieser Theorie vieles sehr vereinfachen, wenn man voraussetzt, dass sich Polygone in Dreiecke zerlegen lassen. Das lässt sich auf folgende Weise beweisen.

Zunächst ist klar, dass sich jede Figur durch eine einzige Teilungslinie (gerade, kürzeste) in zwei Figuren zerlegen lässt, die zusammen einen einspringenden Winkel weniger haben. Durch wiederholte Teilung dieser Art zerlegt man das Polygon also in andere, die keine einspringenden Winkel mehr haben.

Nun sei $P$ ein Polygon mit ausschließlich ausspringenden Winkeln an den Punkten $(1), (2), (3) .\,.\,. (n)$. Ich behaupte, dass die Linie $(1)(3)$ ganz im Inneren des Polygons liegt.
